\chapter[Dự đoán giá ô tô cũ]{Dự đoán giá ô tô cũ ở TP.HCM và Hà Nội}
\setlength{\emergencystretch}{3em}

\section{Bài toán và hướng tiếp cận}

Giá một chiếc ô tô cũ phụ thuộc vào nhiều thứ cùng lúc: dòng xe, đời xe, số km, phiên bản, nơi bán. Người mua khó biết
một mức giá rao là hợp lý hay không, còn người bán khó biết nên rao bao nhiêu. Bài tập này xây dựng mô hình ước lượng
\textbf{giá rao bán} của một chiếc ô tô đã qua sử dụng từ các thông tin có trong tin đăng, đồng thời trả lời ba câu hỏi:

\begin{itemize}
  \item Tin rao tự thu thập thiếu những gì, thiếu vì sao, và nên điền như thế nào để có cơ sở?
  \item Mỗi năm tuổi, một chiếc xe mất bao nhiêu phần trăm giá, và hãng nào giữ giá tốt nhất?
  \item Khi cùng dòng xe và cùng đời, số km, thành phố và xuất xứ ảnh hưởng tới giá đến đâu?
\end{itemize}

Đây là bài toán \textbf{hồi quy có giám sát}. Quy trình gồm: kiểm tra chất lượng dữ liệu, đặc biệt là giá trị thiếu;
giới hạn bài toán vào xe đã qua sử dụng thuộc các dòng phổ biến; chia train / validation / test; phân tích khám phá trên
train; tạo feature, điền giá trị thiếu, mã hóa và chuẩn hóa trong một \texttt{Pipeline} của scikit-learn
\cite{tab-sklearn}; so sánh hai baseline với năm mô hình hồi quy; chấm trên tập test và phân tích lỗi. Notebook chạy
được trên Google Colab với \textit{Runtime > Run all} trong khoảng một phút, không cần GPU.

\section{Bộ dữ liệu}

\begin{itemize}
  \item \textbf{Nguồn:} chuyên mục \textit{Ô tô} của Chợ Tốt \cite{tab-chotot}. Nhóm tự thu thập qua API công khai mà
        trang web dùng để hiển thị danh sách tin, không cần đăng nhập, khoảng một request mỗi giây. Script thu thập là
        \texttt{crawl\_chotot\_oto.py}. Kiểu dáng, số ghế, xuất xứ, màu, hộp số, nhiên liệu được API trả về dạng mã
        số; script tra nhãn chữ của mỗi mã bằng API chi tiết của một tin có mã đó.
  \item \textbf{Phạm vi và thời điểm:} toàn bộ tin đang hiển thị ở TP.HCM và Hà Nội chiều 20/11/2026, gồm cả xe mới và
        xe đã qua sử dụng; tin cũ nhất được đăng ngày 21/09/2026.
  \item \textbf{Ẩn danh:} file CSV không chứa tên người bán, ID tài khoản, ảnh, phần mô tả tự do, địa chỉ, tọa độ và biển
        số; số điện thoại và biển số trong tiêu đề được che. Về người bán chỉ giữ các con số tổng hợp. Dữ liệu chỉ dùng
        cho mục đích học tập.
  \item \textbf{Lưu ý về nhãn:} giá trong dữ liệu là giá rao do người bán tự đặt, không phải giá giao dịch.
\end{itemize}

Bộ dữ liệu có 12~237 tin (8~840 ở TP.HCM, 3~397 ở Hà Nội; 10~686 xe đã dùng, 1~551 xe mới) thuộc 69 hãng và 435 dòng
xe, với 30 cột (Bảng~\ref{tab:tab-cols}), trong đó 13 cột có ô thiếu.

\begin{table}[H]\centering
\caption{Các nhóm cột của bộ dữ liệu.}\label{tab:tab-cols}
\begin{tabular}{>{\raggedright\arraybackslash}p{2.4cm}>{\raggedright\arraybackslash}p{12cm}}
\toprule
Nhóm & Cột \\
\midrule
Tin đăng & \texttt{list\_id}, \texttt{subject}, \texttt{posted\_at}, \texttt{orig\_posted\_at}, \texttt{crawled\_at}, \texttt{number\_of\_images}, \texttt{has\_video}, \texttt{is\_price\_not\_valid} \\
Nơi bán & \texttt{city}, \texttt{district} \\
Chiếc xe & \texttt{brand}, \texttt{model}, \texttt{mfg\_year}, \texttt{body\_type}, \texttt{seats}, \texttt{gearbox}, \texttt{fuel}, \texttt{origin}, \texttt{color}, \texttt{condition}, \texttt{mileage\_km} \\
Giấy tờ & \texttt{owners}, \texttt{has\_registration}, \texttt{has\_accessories} \\
Người bán & \texttt{seller\_is\_shop}, \texttt{seller\_live\_ads}, \texttt{seller\_sold\_ads}, \texttt{seller\_avg\_rating}, \texttt{seller\_total\_rating} \\
Target & \texttt{price\_vnd} (giá rao, VND) \\
\bottomrule
\end{tabular}
\end{table}

\reportfig{fig01_categorical.png}{Phân phối các cột phân loại chính. Ba hãng nhiều tin nhất (Toyota, VinFast, Ford)
chiếm 42\% số tin; 25\% tin không ghi kiểu dáng (thanh xám); 11\% là xe điện và 85\% là xe số tự động.}{fig:tab-cat}

\section{Kiểm tra chất lượng dữ liệu}\label{sec:tab-quality}

\subsection{Giá trị thiếu: ba kiểu thiếu khác nhau}

Phần lớn công sức kiểm tra dành cho giá trị thiếu, vì cách điền chỉ có cơ sở khi biết dữ liệu thiếu theo kiểu nào.
Dữ liệu có ba kiểu thiếu:

\begin{enumerate}
  \item \textbf{Thiếu vì không áp dụng.} Ở tin xe mới, số km, số đời chủ, đăng kiểm và phụ kiện trống gần như 100\% vì
        form đăng xe mới không hỏi các mục này (Hình~\ref{fig:tab-miss}). Bài toán giới hạn ở xe đã dùng nên kiểu thiếu
        này không cần điền.
  \item \textbf{Thiếu ẩn.} 490 tin ghi số ghế \texttt{Khác}; 177 tin xe đã dùng khai 0 km và 479 tin khai dưới 1~000 km
        dù xe đã từ 2 năm tuổi (một chiếc Lexus GX 460 đời 2010 khai 0 km); 55 tin khai trên 500~000 km với các số như
        888~888, 999~999. Các giá trị này được đổi thành ô thiếu.
  \item \textbf{Người bán bỏ trống.} Ở xe đã dùng: \texttt{owners} 50\%, \texttt{has\_accessories} 48\%,
        \texttt{has\_registration} 44\%, \texttt{origin} 39\%, \texttt{body\_type} 25\%, \texttt{seats} 19\% (tính cả
        giá trị \texttt{Khác}), \texttt{color} 6\%. Hãng, dòng xe, năm sản xuất, hộp số, nhiên liệu và giá không thiếu.
\end{enumerate}

\reportfig{fig03_missing_by_condition.png}{Tỉ lệ thiếu theo cột, tách theo tình trạng xe. Màu cam: các cột gần như luôn
trống ở xe mới vì không áp dụng.}{fig:tab-miss}

Để biết kiểu thứ ba có ngẫu nhiên không, nhóm làm ba kiểm tra trên xe đã dùng (Bảng~\ref{tab:tab-mar},
Hình~\ref{fig:tab-misscorr}):

\begin{itemize}
  \item \textbf{Thiếu cùng nhau:} ba cột giấy tờ có cờ ``bị thiếu'' tương quan 0.73 -- 0.90 với nhau: người bán bỏ qua cả
        nhóm mục một lúc. \texttt{body\_type} và \texttt{seats} tương quan 0.40.
  \item \textbf{Ai bỏ trống:} người bán cá nhân bỏ trống nhiều hơn salon ở \texttt{body\_type}, \texttt{owners},
        \texttt{has\_accessories} (chênh 10 -- 11 điểm \%); các cột khác chênh không quá 4 điểm \%.
  \item \textbf{Thiếu và giá:} giá của mỗi tin được chia cho giá median của các tin cùng dòng xe và cùng năm sản xuất. Với
        cả 7 cột, median của tỉ số này ở tin bị thiếu và tin đầy đủ đều bằng 1.00.
\end{itemize}

\begin{table}[H]\centering
\caption{Kiểm tra giá trị thiếu do người bán bỏ trống (xe đã dùng). Giá tương đối: giá chia cho median giá các tin cùng
dòng xe và cùng năm sản xuất (nhóm từ 5 tin).}\label{tab:tab-mar}
\small
\begin{tabular}{lrrrrr}
\toprule
& & \multicolumn{2}{c}{\% thiếu theo người bán} & \multicolumn{2}{c}{Giá tương đối} \\
\cmidrule(lr){3-4}\cmidrule(lr){5-6}
Cột & \% thiếu & salon & cá nhân & khi thiếu & khi có \\
\midrule
\texttt{body\_type} & 24.7 & 21.9 & 32.0 & 1.00 & 1.00 \\
\texttt{seats} & 19.0 & 19.0 & 18.8 & 1.00 & 1.00 \\
\texttt{origin} & 39.5 & 38.5 & 42.2 & 1.00 & 1.00 \\
\texttt{color} & 5.9 & 6.2 & 5.1 & 1.00 & 1.00 \\
\texttt{owners} & 50.4 & 47.6 & 57.9 & 1.00 & 1.00 \\
\texttt{has\_registration} & 44.5 & 44.0 & 45.9 & 1.00 & 1.00 \\
\texttt{has\_accessories} & 48.1 & 45.0 & 56.5 & 1.00 & 1.00 \\
\bottomrule
\end{tabular}
\end{table}

\reportfig{fig04_missing_correlation.png}{Tương quan giữa các cờ ``bị thiếu''. Ba cột giấy tờ thiếu cùng nhau; các cặp
khác tương quan yếu.}{fig:tab-misscorr}

Kết luận: việc thiếu phụ thuộc vào người bán và cách họ điền form, tức vào những thứ quan sát được, chứ không phụ thuộc
vào giá. Vì vậy có thể điền ô thiếu dựa trên các cột khác mà không làm lệch mô hình; cột nào điền được và điền bằng gì
được kiểm tra bằng số ở Mục~\ref{sec:tab-impute}.

\subsection{Các lỗi khác}

Bảng~\ref{tab:tab-issues} tổng kết các vấn đề và cách xử lý. Đáng chú ý là lỗi giá: cờ \texttt{is\_price\_not\_valid}
của Chợ Tốt không đánh dấu tin nào; 32 tin dưới 20 triệu phần lớn là tin dịch vụ, tin tìm mua hoặc giá gõ thiếu số (một
chiếc Mazda CX-5 2015 rao 380~000 đồng); còn những lỗi khó thấy hơn, như tin rao tiền trả trước hoặc giá gấp nhiều lần
xe cùng loại, chỉ lộ ra khi so với xe cùng dòng và đời gần nhau (Mục~\ref{sec:tab-eda}). Salon cũng hay đăng lại cùng
một chiếc xe: 179 tin trùng được bỏ trước khi chia tập để không có cùng một chiếc xe nằm ở cả train và test.

\begin{table}[H]\centering
\caption{Các vấn đề dữ liệu và cách xử lý.}\label{tab:tab-issues}
\begin{tabular}{>{\raggedright\arraybackslash}p{3.6cm}>{\raggedright\arraybackslash}p{6.4cm}>{\raggedright\arraybackslash}p{4.6cm}}
\toprule
Vấn đề & Kết quả kiểm tra & Cách xử lý \\
\midrule
Thiếu vì không áp dụng & xe mới không có số km, số đời chủ, giấy tờ & giới hạn ở xe đã dùng \\
Thiếu ẩn & 490 tin số ghế \texttt{Khác}; 656 tin số km quá nhỏ so với tuổi xe; 55 tin trên 500~000 km & đổi thành ô thiếu \\
Người bán bỏ trống & kiểu dáng, số ghế, xuất xứ, giấy tờ, màu, các cột người bán; không liên quan tới giá & điền theo dòng xe hoặc nhóm \texttt{không rõ}, kèm cột cờ \\
Giá rác & 32 tin dưới 20 triệu & bỏ tin \\
Giá vô lý so với xe cùng loại & tiền trả trước, giá gấp nhiều lần xe gần giống & bỏ theo ngưỡng học từ train \\
Tin trùng & 179 tin đăng lại cùng một chiếc xe & bỏ, giữ tin mới nhất \\
Năm trong tiêu đề khác \texttt{mfg\_year} & 1.3\% lệch hơn 1 năm & giữ \texttt{mfg\_year} \\
Dòng xe hiếm, nhóm \texttt{Dòng khác} & 212 dòng có dưới 5 tin & giữ 85 dòng có từ 30 tin \\
Giá lệch phải & skew trên 4 ở thang gốc & học trên $\log_{10}$(giá) \\
\bottomrule
\end{tabular}
\end{table}

\subsection{Phạm vi bài toán và chia tập}

Bài toán giới hạn ở \textbf{xe đã qua sử dụng}. Sau các luật làm sạch còn 10~479 tin thuộc 434 dòng xe, nhưng 212 dòng
có dưới 5 tin. Nhóm giữ 85 dòng có từ 30 tin, phủ 77.9\% số tin; các nhóm \texttt{Dòng khác} (Chợ Tốt gom các dòng xe
không có trong danh mục) không được giữ dù đủ số tin. Tập còn 8~161 tin, chia 70 / 15 / 15 phân tầng theo 10 khoảng giá:
train 5~712, validation 1~224, test 1~225 tin. Mọi phân tích và thống kê sau đó chỉ dùng train; test chỉ được chấm một
lần.

\section{Phân tích khám phá dữ liệu}\label{sec:tab-eda}

Giá ở thang gốc lệch phải mạnh (skew 4.2, median 495 triệu); sau khi lấy $\log_{10}$, phân phối gần đối xứng (skew
$-0.29$), nên mô hình học trên $\log_{10}$(giá). Riêng dòng xe đã giải thích 58\% phương sai của log giá.

\textbf{Mỗi năm xe mất bao nhiêu giá?} So giá trực tiếp giữa các hãng sẽ bị lẫn với việc mỗi hãng bán những dòng xe khác
nhau, nên với mỗi hãng nhóm ước lượng hồi quy $\log(\text{giá}) \sim \text{tuổi xe} + \text{dòng xe}$ trên xe 1 -- 12
năm tuổi. Tính chung, giá giảm khoảng 8.9\% mỗi năm khi so trong cùng dòng xe; Mazda giữ giá tốt nhất (6.3\%/năm),
Mercedes Benz mất giá nhanh nhất (13.2\%/năm), Toyota ở giữa (8.3\%/năm) (Hình~\ref{fig:tab-dep}).

\reportfig{fig08_depreciation_by_brand.png}{Giá giảm mỗi năm tuổi khi so trong cùng dòng xe, theo hãng (hệ số tuổi xe
của hồi quy log giá theo tuổi xe và dòng xe, xe 1 -- 12 năm tuổi). Đường nét đứt: 8.9\% tính chung.}{fig:tab-dep}

\textbf{Số km, thành phố, xuất xứ.} Khi cố định dòng xe và đời xe, xe chạy dưới 5~000 km/năm được rao 1.02 lần giá
median, xe chạy trên 30~000 km/năm 0.98 lần (Hình~\ref{fig:tab-km}). Giá ở Hà Nội bằng 0.993 lần TP.HCM (486 cặp dòng
xe + đời), xe nhập khẩu bằng 0.995 lần xe lắp ráp trong nước (241 cặp), xe của salon bằng 1.011 lần xe của người bán cá
nhân (509 cặp). Chênh lệch giá thô giữa các nhóm này chủ yếu đến từ việc mỗi nhóm bán những dòng xe và đời xe khác nhau.

\reportfig{fig09_km_effect.png}{Giá so với median của các tin cùng dòng xe và cùng đời, theo số km mỗi năm. Ảnh hưởng
của số km chỉ vài phần trăm khi đã biết dòng xe và đời xe.}{fig:tab-km}

\textbf{Outlier.} Theo quy tắc IQR, 4.6\% giá ở thang gốc là outlier, chỉ còn 3.7\% ở thang log; phần lớn là xe sang
hoặc xe rất cũ có thật. Lỗi giá thật sự lộ ra khi so mỗi tin với các tin train cùng dòng xe, năm sản xuất chênh không quá
2 năm: 7 tin rao gấp hơn 3 lần (một chiếc Ford Ranger 2019 rao 4~355 triệu trong khi tiêu đề ghi 435 triệu) và 31 tin
dưới 30\% giá của những chiếc gần giống.

\section{Tiền xử lý dữ liệu}\label{sec:tab-prep}

\begin{itemize}
  \item \textbf{Bỏ tin có giá vô lý:} giá dưới 30\% hoặc trên 3 lần median giá của các tin train cùng dòng xe, năm sản
        xuất chênh không quá 2 năm. Bỏ 38 tin ở train, 3 ở validation, 9 ở test.
  \item \textbf{Tạo feature:} \texttt{age} (tuổi xe lúc thu thập), \texttt{log\_km}, \texttt{km\_per\_year} và bốn cờ từ
        khóa trong tiêu đề (bản cao cấp, dẫn động 4x4, xe lướt, vay / trả góp).
  \item \textbf{Điền giá trị thiếu} theo nguyên nhân thiếu của từng cột (Mục~\ref{sec:tab-impute}).
  \item \textbf{Mã hóa và chuẩn hóa:} one-hot cho cột phân loại ít mức; target encoding \cite{tab-te} kèm one-hot cho
        hãng, dòng xe và quận; \texttt{StandardScaler} cho cột số. Tất cả nằm trong một \texttt{ColumnTransformer} đặt đầu
        \texttt{Pipeline}, nên mọi thống kê chỉ học từ train.
\end{itemize}

\subsection{Điền giá trị thiếu theo dòng xe}\label{sec:tab-impute}

Kiểu dáng, số ghế và xuất xứ gần như được quyết định bởi dòng xe (Toyota Innova là MPV 8 chỗ lắp ráp trong nước). Nhóm
viết một transformer của scikit-learn, \texttt{GroupModeImputer}: \texttt{fit} học trên train bảng ``dòng xe $\to$ giá
trị phổ biến nhất'', \texttt{transform} điền ô thiếu bằng giá trị của cùng dòng xe, nếu không có thì của cùng hãng, rồi
của toàn bộ train. Mỗi cột được điền còn có thêm một cột cờ ``bị thiếu''.

Cách điền được kiểm tra trực tiếp: trên train, giấu ngẫu nhiên 20\% số ô đang có giá trị của một cột, điền lại rồi so
với giá trị thật, lặp 5 lần (Hình~\ref{fig:tab-impcheck}). Điền theo dòng xe đúng 90\% với \texttt{body\_type}, 91\% với
\texttt{seats} và 67\% với \texttt{origin}, so với 42\%, 53\% và 53\% khi điền một giá trị chung. Với \texttt{color},
\texttt{owners}, \texttt{has\_registration}, điền theo dòng xe không hơn điền chung. Tỉ lệ đúng cao của hai cột giấy tờ
chỉ vì gần như mọi ô có giá trị đều là \texttt{1 chủ} và \texttt{Có}, trong khi ô trống ở đây có thể không ngẫu nhiên
(người bán có xe nhiều đời chủ có lý do để bỏ qua). Vì vậy bốn cột màu và giấy tờ được giữ thành nhóm riêng
\texttt{không rõ}. Các cột số (số km, các cột người bán) được điền bằng median của train kèm cột cờ.

\reportfig{fig12_imputation_check.png}{Điền lại các ô bị giấu trên train: tỉ lệ điền đúng giá trị thật khi điền một giá
trị phổ biến nhất chung (xám) và giá trị phổ biến nhất của cùng dòng xe (xanh). Trung bình của 5 lần giấu.}{fig:tab-impcheck}

\section{Xây dựng mô hình}

Thước đo chính là $R^2$ trên $\log_{10}$(giá), kèm sai số tuyệt đối theo \% giá thật lấy trung vị trên các tin
(\textit{sai số \% median}), tỉ lệ tin đoán lệch không quá 10\% và MAE (triệu đồng). Mỗi mô hình được chấm bằng
cross-validation 5 fold trên train rồi trên validation. Hai mô hình chỉ được coi là khác nhau khi hiệu $R^2$ validation
lớn hơn độ lệch giữa các fold; nếu nhỏ hơn, chọn mô hình có sai số \% median thấp hơn.

Hai baseline là đoán giá trung bình của train và đoán giá median của các xe cùng dòng và cùng năm sản xuất (cách một
người mua thường tra giá). Năm mô hình dùng chung bộ tiền xử lý: hồi quy tuyến tính, kNN, cây quyết định, random forest
\cite{tab-rf} và gradient boosting (\texttt{HistGradientBoostingRegressor}). Bảng~\ref{tab:tab-val} là kết quả trên
validation.

\begin{table}[H]\centering
\caption{Kết quả trên validation (1~221 tin sau khi bỏ giá vô lý).}\label{tab:tab-val}
\small
\begin{tabular}{lrrr}
\toprule
Mô hình & $R^2$ CV (± độ lệch) & $R^2$ validation & Sai số \% median \\
\midrule
Baseline: giá trung bình & $-0.001$ ± 0.001 & 0.000 & 36.3\% \\
Baseline: median cùng dòng xe và đời & 0.858 ± 0.009 & 0.820 & 6.6\% \\
Hồi quy tuyến tính & 0.928 ± 0.004 & 0.923 & 8.0\% \\
kNN ($k = 10$) & 0.791 ± 0.012 & 0.791 & 14.6\% \\
Cây quyết định (sâu 12) & 0.919 ± 0.007 & 0.924 & 7.8\% \\
\textbf{Random forest} & 0.946 ± 0.004 & \textbf{0.953} & \textbf{5.6\%} \\
Gradient boosting & 0.949 ± 0.004 & 0.953 & 6.0\% \\
\bottomrule
\end{tabular}
\end{table}

Random forest và gradient boosting cùng có $R^2$ validation 0.953, chênh nhau ít hơn độ lệch giữa các fold (0.004);
random forest được chọn vì sai số \% median thấp hơn (5.6\% so với 6.0\%). Baseline median cùng dòng xe và đời đã khá
mạnh (sai số 6.6\%): biết dòng xe và năm sản xuất là đoán được gần đúng giá. kNN kém nhất vì khoảng cách được tính trên
hơn 200 cột sau one-hot. Khi thay đổi tham số, cây quyết định có khoảng cách giữa $R^2$ train và validation tăng từ
0.008 (sâu 4) lên 0.043 (không giới hạn độ sâu), dấu hiệu học thuộc train; kNN tốt nhất ở $k = 5$; gradient boosting
gần như đi ngang sau khoảng 400 vòng lặp.

\subsection{Cách điền giá trị thiếu ảnh hưởng tới mô hình ra sao?}\label{sec:tab-impeffect}

Random forest được huấn luyện theo ba cách xử lý \texttt{body\_type}, \texttt{seats}, \texttt{origin}: bỏ khỏi train mọi
tin có ô thiếu; điền đơn giản (median cho số ghế, nhóm \texttt{không rõ} cho hai cột còn lại); điền theo dòng xe
(Bảng~\ref{tab:tab-impeffect}, Hình~\ref{fig:tab-impeffect}).

\begin{table}[H]\centering
\caption{Ba cách xử lý giá trị thiếu (random forest, chấm trên validation).}\label{tab:tab-impeffect}
\begin{tabular}{lrrrrr}
\toprule
& & \multicolumn{2}{c}{Toàn bộ validation} & \multicolumn{2}{c}{612 tin có ô thiếu} \\
\cmidrule(lr){3-4}\cmidrule(lr){5-6}
Cách & Tin train & $R^2$ & Sai số \% & $R^2$ & Sai số \% \\
\midrule
Bỏ tin có ô thiếu & 2~782 & 0.939 & 6.7\% & 0.935 & 7.6\% \\
Điền đơn giản & 5~674 & 0.953 & 5.7\% & 0.954 & 6.3\% \\
\textbf{Điền theo dòng xe} & 5~674 & \textbf{0.953} & \textbf{5.6\%} & \textbf{0.955} & \textbf{5.8\%} \\
\bottomrule
\end{tabular}
\end{table}

\reportfig{fig15_imputation_effect.png}{Sai số \% median trên validation theo cách xử lý giá trị thiếu: toàn bộ
validation (trái) và 612 tin có ít nhất một ô thiếu (phải).}{fig:tab-impeffect}

Bỏ tin có ô thiếu làm mất 51\% train và cho sai số cao nhất. Hai cách điền gần như ngang nhau trên toàn bộ validation,
vì dòng xe đã có sẵn trong mô hình qua \texttt{brand\_model}; khác biệt nằm ở các tin có ô thiếu, nơi điền theo dòng xe
cho sai số 5.8\% so với 6.3\%. Nhóm giữ cách điền theo dòng xe: đúng hơn khi kiểm tra từng ô và tốt hơn ở các tin có ô
thiếu.

\section{Đánh giá trên tập test}

Mọi mô hình được chấm một lần trên test (Bảng~\ref{tab:tab-test}). Random forest đạt $R^2$ 0.953, bằng validation
(0.953) và gần cross-validation (0.946); một nửa số xe được đoán lệch không quá 5.5\% giá rao, so với 7.2\% nếu chỉ tra
giá median của xe cùng dòng và đời. Gradient boosting có $R^2$ gần như bằng (0.952) nhưng sai số \% median cao hơn
(6.3\%), nhất quán với kết luận ở validation.

\begin{table}[H]\centering
\caption{Kết quả trên test (1~216 tin), $R^2$ tính trên $\log_{10}$(giá). Hàng in đậm là mô hình cuối cùng, chọn theo
validation.}\label{tab:tab-test}
\small
\begin{tabular}{lrrrr}
\toprule
Mô hình & $R^2$ & Sai số \% median & Lệch $\le$ 10\% & MAE (triệu) \\
\midrule
Baseline: giá trung bình & 0.000 & 36.0\% & 14.2\% & 275.7 \\
Baseline: median cùng dòng xe và đời & 0.847 & 7.2\% & 60.4\% & 74.7 \\
Hồi quy tuyến tính & 0.932 & 7.9\% & 58.7\% & 72.4 \\
kNN ($k = 10$) & 0.812 & 14.6\% & 36.8\% & 122.6 \\
Cây quyết định (sâu 12) & 0.917 & 8.0\% & 56.9\% & 72.6 \\
\textbf{Random forest} & \textbf{0.953} & \textbf{5.5\%} & \textbf{70.3\%} & \textbf{53.9} \\
Gradient boosting & 0.952 & 6.3\% & 66.9\% & 58.4 \\
\bottomrule
\end{tabular}
\end{table}

\reportfig{fig17_pred_vs_true.png}{Giá dự đoán so với giá rao thật trên test (thang log). Dải xám: lệch không quá 20\%,
chứa phần lớn số tin; các điểm của Hà Nội và TP.HCM nằm lẫn vào nhau.}{fig:tab-pred}

\reportfig{fig18_feature_importance.png}{Permutation importance trên test: $R^2$ giảm bao nhiêu khi xáo trộn từng cột đầu
vào. Dòng xe (0.636) và tuổi xe (0.626) quan trọng hơn hẳn mọi cột khác.}{fig:tab-imp}

\subsection{Mô hình sai ở đâu?}

Sai số cao nhất ở 20\% xe rẻ nhất (dưới 305 triệu): 10.3\%, và mô hình thường đoán cao hơn giá rao (tỉ số dự đoán / thật
1.04); ở 20\% xe đắt nhất mô hình hơi đoán thấp (0.97). Tin của salon dễ đoán hơn tin của người bán cá nhân ở cả hai
thành phố (5.2\% và 5.5\% so với 6.7\% và 6.5\%). Các tin sai nhiều nhất là xe đời cũ giá thấp hoặc tin mà giá không
phải giá xe, ví dụ một tin Hyundai Accent ghi ``đưa mình 185tr ra góp tiếp'', tức giá là phần tiền góp còn lại.

\section{Dự đoán giá một chiếc xe}

Hàm \texttt{predict\_price()} nhận thông tin như khi đăng một tin trên Chợ Tốt và đi qua đúng pipeline của mô hình cuối
cùng; xe thuộc dòng ngoài phạm vi bị từ chối. Bảng~\ref{tab:tab-pricing} thay đổi lần lượt từng thông tin của một chiếc
Toyota Innova.

\begin{table}[H]\centering
\caption{Giá dự đoán khi thay đổi từng thông tin của một chiếc Toyota Innova 2019, 80~000 km, người bán cá nhân ở Quận
Gò Vấp.}\label{tab:tab-pricing}
\begin{tabular}{lr}
\toprule
Thay đổi & Giá dự đoán (triệu) \\
\midrule
Xe gốc & 373.2 \\
Đời 2015 thay vì 2019 & 252.0 \\
Đi 150~000 km thay vì 80~000 km & 375.9 \\
Bán ở Hà Nội (Cầu Giấy) & 367.8 \\
Do salon bán & 375.9 \\
Không khai số ghế, kiểu dáng, xuất xứ & 373.0 \\
Toyota Fortuner thay vì Innova & 561.3 \\
\bottomrule
\end{tabular}
\end{table}

Lùi đời về 2015 làm giá giảm 32\%, khoảng 9\% mỗi năm, khớp với Hình~\ref{fig:tab-dep}. Số km, thành phố và loại người
bán chỉ làm giá chênh khoảng 1\%, khớp với Hình~\ref{fig:tab-km}. Bỏ trống ba cột được điền theo dòng xe cho kết quả gần
như y hệt, vì pipeline điền lại đúng các giá trị của Innova. Đổi sang Fortuner thì giá lên 561 triệu.

\section{Kết luận}

Trên 85 dòng xe phổ biến (78\% số tin xe đã dùng sau khi làm sạch), random forest giải thích 95.3\% biến động của log
giá trên tập test và đoán một nửa số xe lệch không quá 5.5\% giá rao, so với 7.2\% của cách tra giá median cùng dòng
xe và đời. Giá trị thiếu trong tin rao tự thu thập có ba kiểu (không áp dụng, thiếu ẩn, người bán bỏ trống); kiểu thứ
ba phụ thuộc vào người bán chứ không liên quan tới giá, nên có thể điền dựa trên các cột khác. Điền theo dòng xe đúng
90 -- 91\% với kiểu dáng và số ghế, và cho sai số thấp nhất trên các tin có ô thiếu. Xe mất trung bình 8.9\% giá mỗi
năm tuổi, từ 6.3\% (Mazda) tới 13.2\% (Mercedes Benz); số km, thành phố và xuất xứ chỉ làm giá chênh vài phần trăm khi
đã biết dòng xe và đời xe.

\subsection*{Hạn chế}
\begin{itemize}
  \item \textbf{Giá rao không phải giá bán}, và dữ liệu chỉ là một lần thu thập (20/11/2026): xe bán nhanh có thể đã bị
        gỡ trước đó.
  \item \textbf{Thiếu thông tin tình trạng xe.} Đâm đụng, ngập nước, bảo dưỡng nằm trong phần mô tả và ảnh mà nhóm không
        thu thập để bảo vệ thông tin cá nhân; phiên bản xe chỉ được nhận ra qua vài từ khóa trong tiêu đề.
  \item \textbf{Một nửa số tin không có thông tin giấy tờ}, nên mô hình gần như không học được ảnh hưởng của số đời chủ
        và đăng kiểm.
  \item \textbf{Chỉ áp dụng cho 85 dòng xe phổ biến}; dòng hiếm và các nhóm \texttt{Dòng khác} nằm ngoài phạm vi.
  \item \textbf{Luật bỏ giá vô lý dùng nhãn:} 12 tin ở validation và test bị bỏ dựa trên chính giá rao của chúng.
  \item \textbf{Năm sản xuất do người bán khai}; 1.3\% tiêu đề ghi năm lệch hơn 1 năm so với cột này.
\end{itemize}

\subsection*{Hướng mở rộng}
Thu thập định kỳ để theo dõi giá theo thời gian; lấy phiên bản xe có cấu trúc từ trang chi tiết của tin thay cho từ khóa
trong tiêu đề; trích tình trạng xe từ phần mô tả sau khi ẩn danh; dự đoán một khoảng giá (quantile regression) thay vì
một con số.

\begin{thebibliography}{9}
\bibitem{tab-chotot} Chợ Tốt. \textit{Mua bán ô tô cũ}. \url{https://xe.chotot.com/mua-ban-oto}, truy cập ngày
  20/11/2026.
\bibitem{tab-sklearn} F. Pedregosa và cộng sự. \textit{Scikit-learn: Machine Learning in Python}. Journal of Machine
  Learning Research, 12, 2011, tr. 2825--2830.
\bibitem{tab-te} D. Micci-Barreca. \textit{A Preprocessing Scheme for High-Cardinality Categorical Attributes in
  Classification and Prediction Problems}. ACM SIGKDD Explorations Newsletter, 3(1), 2001, tr. 27--32.
\bibitem{tab-rf} L. Breiman. \textit{Random Forests}. Machine Learning, 45(1), 2001, tr. 5--32.
\end{thebibliography}
